\section{What a finite-resource zero estimate means}
\label{sec:operational}

The empirical distinction between a central residual and a reported
interpolation error motivates a definition based on a complete measurement
protocol. Nothing in this observational definition makes the mathematical
ordinate $\gamma_j$ evolve with cosmic time.

A protocol $\mathcal P$ specifies an encoding, a measured signal, a likelihood
$p(y\mid\theta,\boldsymbol\eta,\mathcal P)$, an estimator, and resources
such as interrogation duration, repetitions, bandwidth, and calibration.
Here $\theta$ is an unknown signal coordinate and $\boldsymbol\eta$ denotes
nuisance parameters. The ideal encoding places the relevant signal feature
at $\theta=\gamma_j$. For a tolerance $\Delta_j>0$ and a chosen error
probability $\epsilon_{\rm conf}$, a pointwise operational requirement is
\begin{equation}
 \Pr_{\mathcal P}\!\left(
   |\widehat\gamma_j-\gamma_j|\leq\Delta_j
 \right)\geq 1-\epsilon_{\rm conf}.
 \label{eq:accessible_criterion}
\end{equation}
Coverage must be justified over the declared nuisance-parameter range, not
only at a fitted calibration value. The tolerance may be absolute or relative
to a specified neighboring gap, but the choice must precede the comparison.
An uncertainty bar without its coverage and systematic-error model does not
establish Eq.~\eqref{eq:accessible_criterion}.

Inference in this definition uses the measured response. Returning a known
ordinate from a table is excluded as an independent analog measurement.
For a validation experiment, blinded signal coordinates or held-out target
windows can prevent reference information from being counted as experimental
precision. Engineered responses remain useful tests of control and encoding,
even when their construction already uses a mathematical function.

For a resource budget $\mathbf R$, let $\mathcal A(\mathbf R)$ contain the
targets meeting Eq.~\eqref{eq:accessible_criterion} in the declared protocol
class. An isolated-target height is
$T_{\rm target}=\sup\{\gamma_j:j\in\mathcal A(\mathbf R)\}$.
A complete-prefix requirement is stronger: one jointly budgeted procedure
must estimate every target through $J$, with simultaneous coverage
\begin{equation}
 \Pr\!\left(\bigcap_{j=1}^{J}
  \{|\widehat\gamma_j-\gamma_j|\leq\Delta_j\}\right)
 \geq1-\epsilon_{\rm conf}.
 \label{eq:prefix_coverage}
\end{equation}
The corresponding height is $\gamma_J$. Pointwise coverage does not imply
this joint guarantee, and isolated successful estimates need not form a
contiguous prefix. A scan ending at its largest reported target therefore
does not measure a fundamental maximum height.

\subsection{An illustrative quantum precision calculation}

Consider the frequency encoding $\omega=\kappa\theta$, where $\kappa$ has
units of inverse time. An initially equatorial qubit evolves under
\begin{equation}
 H_\theta=\frac{\hbar\kappa\theta}{2}\sigma_z,
 \qquad U_\theta(\tau)=\exp(-i\kappa\theta\tau\sigma_z/2).
\end{equation}
The quantum Fisher information per ideal independent interrogation is
$F_Q=(\kappa\tau)^2$. Under the regularity and local unbiasedness assumptions
of the quantum Cram\'er--Rao framework, $M$ independent trials obey
\begin{equation}
 \operatorname{Var}(\widehat\theta)
 \geq\frac{1}{M F_Q}
 =\frac{1}{M\kappa^2\tau^2}.
 \label{eq:ramsey_qfi}
\end{equation}
This is an elementary example within quantum estimation theory
\cite{BraunsteinCaves1994,Giovannetti2006}, not a universal noise model for
every Riemann-zero implementation. For a binary Ramsey signal
$p_+=[1+V\cos(\kappa\theta\tau+\phi)]/2$ with known, parameter-independent
visibility $V$, operation at quadrature gives classical Fisher information
$F=V^2(\kappa\tau)^2$ per shot. Its local statistical bound becomes
$\sigma_\theta\geq1/(V\kappa\tau\sqrt M)$.

A lower bound is not an achieved uncertainty or a coverage guarantee.
Bias, uncertain visibility, correlated shots, and calibration require their
own model. Moreover, the phase is periodic: local Fisher information does
not resolve global aliases without additional interrogation times or prior
range information. A finite-window Fourier width of order $2\pi/\tau$
likewise does not establish a universal prohibition on more precise parameter
estimation; protocols using quantum sensing or external synchronization can
change the relevant resolution behavior \cite{Gefen2019,Boss2017}.
Finite precision follows here from a specified finite-information
experiment under specified assumptions; it does not prove a nonzero error
floor shared by every physically possible protocol.

\subsection{Mean density as a conditional illustration}

The smooth Riemann--von Mangoldt count is
\begin{equation}
 \overline N(T)=\frac{T}{2\pi}\log\frac{T}{2\pi}
 -\frac{T}{2\pi}+\frac78,
 \qquad
 \overline d(T)\simeq\frac{2\pi}{\log(T/2\pi)}.
 \label{eq:mean_spacing}
\end{equation}
The second expression describes the asymptotic mean spacing, not the smallest
individual gap. If a working model supplies an achieved, dimensionless
resolution scale $\delta$ and uses the illustrative criterion
$\delta\leq q\overline d(T)$, with prescribed $0<q<1$, its crossover is
\begin{equation}
 T_{\rm env}=2\pi\exp\!\left(\frac{2\pi q}{\delta}\right).
 \label{eq:mean_envelope}
\end{equation}
This is a mean-density illustration in the asymptotic regime. It is neither
an apparatus-independent bound nor a statement that all lower zeros are
resolvable. It assumes a usable resolution scale rather than substituting a
Cram\'er--Rao lower bound as if it were attained. Bandwidth, signal amplitude,
individual gaps, and calibration may restrict the experiment first. A
conversion from an RMS error to a confidence tolerance additionally requires
a probability model. The exponential sensitivity of
Eq.~\eqref{eq:mean_envelope} makes an absolute height forecast especially
unstable without these ingredients.

\section{A conditional cosmic-time extension}
\label{sec:cosmic_extension}

\subsection{Motivation from the published dynamical model}

The recently published numerical study of a non-autonomous quadratic map
\cite{Wang2026Spectral} investigates a schedule of the form
\begin{equation}
 x_{n+1}=1-\lambda_n x_n^2,
 \qquad
 \lambda_n=\mu_c+\frac{k_1}{\log^2n}
                  +\frac{k_2}{\log^3n},\qquad n>1.
 \label{eq:published_map}
\end{equation}
Its leading inverse-log-squared term is a heuristic numerical ansatz. The
motivation combines the known logarithmic spectral density with an assumed
square-root response of an eigenphase spacing to a structural perturbation.
The resulting transfer-operator correspondence is not a demonstrated
isomorphism or a self-adjoint Hilbert--P\'olya Hamiltonian. The paper reports
larger out-of-sample errors and local spacing statistics inconsistent with
GUE, so its global spectral correspondence is not evidence of matching every
statistical property of the zeta zeros.

Two further assumptions are required to use this motivation here: a mapping
from iteration or spectral scale to cosmic age, and a response law connecting
the evolving structural parameter to a measured precision scale. Neither is
supplied by the numerical correspondence. In particular, preserving the
power two in a noise amplitude assumes a suitable response relation; it
does not follow solely from copying the map schedule.

We therefore introduce the separately testable candidate
\begin{equation}
 \delta_{\rm c}(t)=\delta_{{\rm c}0}
 \left[\frac{\log(t_0/t_*)}{\log(t/t_*)}\right]^p,
 \qquad t>t_*,\quad \delta_{{\rm c}0}\geq0,
 \label{eq:cosmic_ansatz}
\end{equation}
with $p=2$ as the primary hypothesis. Here $\delta_{\rm c}$ is a proposed
contribution to the standard deviation of an encoded ordinate estimator,
not a change in the mathematical ordinate. The scale $t_*$ makes the logarithm
dimensionless, and $t_0$ is a normalization epoch. The amplitude and transfer
to a particular apparatus are not fixed by the earlier map coefficients.
The case $\delta_{{\rm c}0}=0$ is a permitted null model. Neither the
singularity at $t=t_*$ nor an extrapolation into that regime is claimed to
describe early-universe physics.

For a fixed estimator and protocol, one possible phenomenological error
model is
\begin{equation}
 \operatorname{MSE}(\widehat\gamma_j)
 =b_j^2+\sigma_{{\rm stat},j}^2+\sigma_{{\rm cal},j}^2
  +\delta_{{\rm c},j}^2(t).
 \label{eq:conditional_budget}
\end{equation}
This decomposition assumes that the three random error components have zero
mean and are mutually uncorrelated; $b_j$ is the remaining bias. Correlations
require covariance terms. Calibration uncertainty is not automatically
independent random noise, and $\delta_{{\rm c},j}$ may differ across encodings.
Calling it a floor further assumes that the specified component survives
the averaging and control operations allowed by the protocol. Independent
per-shot noise instead decreases on repetition; a correlation model is
needed before asserting an irreducible floor.

For $\delta_{{\rm c}0}>0$, Eq.~\eqref{eq:cosmic_ansatz} gives
\begin{equation}
 \frac{d\log\delta_{\rm c}}{dt}
 =-\frac{p}{t\log(t/t_*)}.
 \label{eq:cosmic_derivative}
\end{equation}
As a scale illustration, choose $t_0=13.8\,\mathrm{Gyr}$ and
$t_*=5.391247\times10^{-44}\,\mathrm{s}$, a conventional Planck-time-scale
normalization rather than an inferred parameter. Using a Julian year,
$\log(t_0/t_*)=140.24423$. For $p=2$, the present fractional rate is
$-1.0334\times10^{-12}\,\mathrm{yr}^{-1}$, corresponding to approximately
$-5.17\times10^{-13}$ in half a year. This is the fractional change of the
hypothesized component, not an observed drift or the fractional change of
the total experimental uncertainty. Its amplitude remains unspecified.

\subsection{Resource growth, identifiability, and falsification}

Let $\mathfrak P(t)$ be the feasible set of protocols including their
calibrated performance at age $t$. If
$\mathfrak P(t_1)\subseteq\mathfrak P(t_2)$ for $t_2>t_1$, with the same
target and confidence criterion, optimization over these sets makes either
accessible-height definition nondecreasing. This is a conditional consequence
of set inclusion. Cosmic age alone does not establish the inclusion:
usable energy, causal access, coherence, and calibration stability need not
improve monotonically. Even in Eq.~\eqref{eq:conditional_budget}, a decreasing
cosmic component does not improve the total error if other components grow.

Entropy and energy limits motivate explicit resource accounting
\cite{Bekenstein1981,MargolusLevitin1998,Lloyd2002}. Converting them into a
zero-height limit would additionally require an encoding, a task definition,
and a justified computational or measurement cost. An explicit table of $J$
values at $b$ bits per value requires roughly $Jb$ bits; an algorithm plus an
index may be far shorter. Finite storage therefore does not by itself select
a largest zero index. The rigorous interval-arithmetic verification of RH
through ordinate $3\times10^{12}$ \cite{PlattTrudgian2021} is a separate
category of evidence from analog spectroscopy, but rules out interpreting
an ordinate of order $4200$ as a general limit on physically computable
zero information.

The available experiments provide no controlled comparison of the same
calibrated apparatus at substantially different cosmic ages. Within each
dataset, changing the zero index or a drive parameter does not vary cosmic
age. A single-epoch residual budget generally cannot identify
$\delta_{{\rm c}0}$ separately from unknown calibration and model discrepancy,
and it cannot select $p=2$ over a constant or another slow time dependence.
Publication dates do not constitute an aging series.

A discriminating test would first fix the encoding, estimator, nuisance
model, and allowed averaging operations, then check the resource dependence
on held-out targets. Repeated independently calibrated measurements could
test the signed drift in Eq.~\eqref{eq:cosmic_derivative}, provided a
nonzero component and sufficient sensitivity were established. A specified
positive-amplitude model could be excluded by a sufficiently stringent
upper limit or by incompatible drift or resource scaling. The unrestricted
family allowing arbitrarily small amplitude cannot be ruled out by a finite
null result. Poor estimates, improved estimates, and null drift must not all
be treated as confirmation of cosmic aging.
